\documentclass[11pt, a4paper]{article}
\usepackage[T1]{fontenc}
\usepackage[utf8]{inputenc}
\usepackage[a4paper, top=2.5cm, bottom=2.5cm, left=2.5cm, right=2.5cm]{geometry}
\usepackage{booktabs}
\usepackage{tabularx}
\usepackage{enumitem}
\usepackage[hidelinks]{hyperref}
\usepackage{parskip}
\usepackage{graphicx}
\usepackage{amsmath}
\usepackage{float}
\setlength{\parindent}{0pt}
\pagestyle{plain}

\begin{document}

% ------------------------------------------------------------------------------
% PORTADA
% ------------------------------------------------------------------------------
\begin{center}
{\Large \textbf{Pontificia Universidad Javeriana}}\\[0.3cm]
Facultad de Ciencias Economicas y Administrativas\\[0.3cm]
Estadistica -- 2026-II\\[1cm]
{\LARGE \textbf{Educacion de los padres y movilidad: ?`pueden los estudiantes de primera generacion cerrar la brecha?}}\\[0.5cm]
{\large Informe Final del Proyecto}\\[1cm]
\textbf{Grupo 6}\\[0.3cm]
Duque Rivera Natalia\\
Flechas Gonzalez Karol Mariana\\
Medrano Marquez Samuel Fernando\\
Salcedo Gonzalez Maria Fernanda\\
Amaya Lara Samuel\\[1cm]
27 de noviembre de 2026
\end{center}

\newpage
\tableofcontents
\newpage

% ------------------------------------------------------------------------------
\section{Introduccion}
% ------------------------------------------------------------------------------
% Contexto: movilidad social intergeneracional en Colombia.
% Pregunta: los estudiantes de primera generacion (padres sin educacion superior)
% obtienen resultados diferentes en SBPro controlando por habilidad previa?
% Concepto de capital cultural y su efecto en la trayectoria educativa.

% ------------------------------------------------------------------------------
\section{Datos y Metodologia}
% ------------------------------------------------------------------------------
% Fuente: cruce Saber 11 -- Saber Pro (ICFES).
% Archivos: cruce\_saber11\_saberpro.parquet, saber11.parquet, saber\_pro.parquet.
% Construccion de la variable primera\_gen: ambos padres SIN educacion superior.
% Variable dependiente 1: PUNT\_GLOBAL\_sbpro.
% Variable dependiente 2: MOD\_INGLES\_PUNT.
% Predictores: PUNT\_GLOBAL\_sb11 (o PUNT\_INGLES\_sb11), primera\_gen,
%   FAMI\_ESTRATOVIVIENDA\_sb11, INST\_ORIGEN\_sbpro.
% Modelo 1 (aditivo), Modelo 2 (interaccion SB11 * primera\_gen),
%   Modelo 3 (ingles con interaccion), Modelo 4 (subsample econ/admin).
% Justificacion: el capital cultural importa mas para ingles?

% ------------------------------------------------------------------------------
\section{Analisis Exploratorio (Fase 1)}
% ------------------------------------------------------------------------------
% Proporcion de estudiantes de primera generacion.
% Distribucion de PUNT\_GLOBAL\_sbpro por primera\_gen.
% Scatter SB11 vs SBPro coloreado por primera generacion.

% ------------------------------------------------------------------------------
\section{Intervalos de Confianza y Pruebas de Hipotesis (Fase 2)}
% ------------------------------------------------------------------------------
% IC para la media de PUNT\_GLOBAL\_sbpro por primera\_gen.
% Prueba t para diferencia de medias primera gen vs no primera gen.
% Prueba de Wilcoxon como alternativa no parametrica.
% Prueba chi-cuadrado para asociacion primera\_gen e INST\_ORIGEN\_sbpro.

% ------------------------------------------------------------------------------
\section{Modelo de Regresion (Fase 3)}
% ------------------------------------------------------------------------------
% Modelo 1: PUNT\_GLOBAL\_sbpro $\sim$ SB11 + primera\_gen + estrato + inst\_origen.
% Modelo 2: PUNT\_GLOBAL\_sbpro $\sim$ SB11 * primera\_gen + estrato.
% Modelo 3: MOD\_INGLES\_PUNT $\sim$ PUNT\_INGLES\_sb11 * primera\_gen + estrato.
% Modelo 4: PUNT\_GLOBAL\_sbpro $\sim$ SB11 * primera\_gen + inst\_origen (econ/admin).
% Tabla de coeficientes con IC al 95\%.
% Comparacion de modelos: R2 ajustado, AIC, BIC.
% Predicciones: primera gen vs no primera gen al P50 de SB11.

% ------------------------------------------------------------------------------
\section{Diagnosticos y Robustez}
% ------------------------------------------------------------------------------
% Diagnosticos para cada modelo: residuales vs ajustados, QQ-plot.
% Test de Breusch-Pagan (homocedasticidad).
% VIF (multicolinealidad).
% Distancia de Cook (observaciones influyentes).
% Errores robustos HC1 para los 4 modelos.

% ------------------------------------------------------------------------------
\section{Resultados y Discusion}
% ------------------------------------------------------------------------------
% La interaccion SB11 * primera\_gen es significativa?
% El capital cultural importa mas para ingles que para puntaje global?
% Graficos de interaccion, coeficientes, predicciones.
% Reflexion sobre movilidad intergeneracional y el rol de la universidad.

% ------------------------------------------------------------------------------
\section{Conclusiones}
% ------------------------------------------------------------------------------
% Respuesta a la pregunta de investigacion.
% Limitaciones del estudio.
% Implicaciones para politica de acompanamiento a primera generacion.

% ------------------------------------------------------------------------------
\section*{Referencias}
% ------------------------------------------------------------------------------
% ICFES, Bourdieu, literatura sobre capital cultural y primera generacion.

% ------------------------------------------------------------------------------
\section*{Anexo: Codigo R}
% ------------------------------------------------------------------------------
% Incluir codigo R utilizado (fase3\_grupo6.R) o referenciarlo.

\end{document}
